\documentclass[12pt]{article}
\usepackage[utf8]{inputenc}
\usepackage[T1]{fontenc}
\usepackage{hyperref}
\title{Ayizi Language Resources for K-12 Curriculum Development}
\author{Compiled by Claude Sonnet 5.5 (claude-sonnet-5-5) for Andr\'as Kornai}
\date{}
\begin{document}
\maketitle

\section{Introduction}
Ayizi (ISO 639-3: yyz; Glottocode ayiz1244) is an endangered Sino-Tibetan (Ngwi, Northern Loloish) language of Yunnan, China. Search results suggest only a few elderly speakers remain. The Ayizi are officially counted as part of the Yi nationality. Their largest village is reported to be Aimailongcun in Shilin County. Some sources say the speakers call themselves Ge, and the language may be related to Gepo, but this is uncertain. The search found no teaching materials.

\section{Primary Educational Resources}
No primers, textbooks or school programmes for Ayizi turned up in the search.

\section{Linguistic and Reference Materials}
\begin{itemize}
\item Related language article (Gepo): \url{https://en.wikipedia.org/wiki/Gepo_language}
\item Endangered Languages Project bibliography page: \url{https://endangeredlanguages.com/lang/7250/bibliography}
\end{itemize}
No grammar or dictionary of Ayizi appeared in the search results.

\section{Religious and Community Resources}
The following people-group profiles were found. They come from mission and advocacy organisations and are not academic sources. I found no scripture translation in Ayizi.
\begin{itemize}
\item Joshua Project profile: \url{https://joshuaproject.net/people_groups/18406/CH}
\item Asia Harvest profile: \url{https://people-groups.asiaharvest.org/China/chinaPeoples/A/Ayizi.pdf}
\end{itemize}

\section{Digital and Audio Resources}
\begin{itemize}
\item Verbix endangered-languages entry: \url{https://docs.verbix.com/EndangeredLanguages/Ayizi}
\end{itemize}
This is a catalogue entry only. No audio or corpus was found.

\section{Curriculum Development Context}
The language is thought to be used by the elderly only, and it is not known to be taught in schools. Documentation appears very thin. A curriculum would depend on first documenting the language with the remaining speakers. Any school use would take place within the Chinese education system, where the Yi nationality is the official category.
\end{document}